\section{Summary}\label{summary}

We present three R packages that address distinct but important gaps in
the econometric toolkit: \texttt{garchur} for unit root testing under
GARCH errors with endogenous structural breaks, \texttt{rbfmvar} for
Residual-Based Fully Modified VAR estimation accommodating mixed
integration orders, and \texttt{tptest} for testing U-shaped and inverse
U-shaped relationships in regression models. These packages implement
methodologies that are widely cited in applied research but previously
lacked dedicated R implementations.

\section{Statement of Need}\label{statement-of-need}

Applied econometric research frequently encounters three challenges that
standard tools handle inadequately:

\begin{enumerate}
\def\labelenumi{\arabic{enumi}.}
\item
  \textbf{Conditional heteroskedasticity in unit root testing}:
  Financial and energy time series often exhibit GARCH-type volatility
  clustering. Standard ADF tests assume homoskedastic errors, leading to
  size distortions and power losses. The Narayan and Liu (2015)
  trend-GARCH unit root test with endogenous structural breaks addresses
  both issues simultaneously, but no R package implements it.
\item
  \textbf{VAR estimation with unknown integration orders}: When a VAR
  system contains an unknown mixture of I(0), I(1), and I(2) variables,
  standard OLS-based VAR estimation produces biased inference. The
  Residual-Based Fully Modified VAR (RBFM-VAR) of Chang (2000),
  extending the Phillips (1995) FM-VAR, applies nonparametric
  corrections that yield valid inference regardless of integration
  orders. This method is particularly valuable when pre-testing for unit
  roots is inconclusive.
\item
  \textbf{Testing for nonlinear (U-shaped) relationships}: The
  Environmental Kuznets Curve, the Laffer curve, and numerous other
  economic theories predict quadratic or other nonlinear functional
  forms with turning points. While ad hoc tests of quadratic coefficient
  significance are common, the formal Lind and Mehlum (2010) test
  provides proper inference on whether a true extreme point exists
  within the data range. No comprehensive R package implements this test
  with extensions to cubic, log-quadratic, and inverse forms.
\end{enumerate}

\section{Packages}\label{packages}

\subsection{garchur}\label{garchur}

Implements the trend-GARCH(1,1) unit root test with endogenous
structural breaks proposed by Narayan and Liu (2015). The procedure
estimates a mean equation with a linear trend and level-shift dummies at
endogenously determined break dates while allowing for GARCH(1,1)
conditional heteroskedasticity. Supports up to three breaks with
critical values interpolated from Monte Carlo simulation tables.

\begin{Shaded}
\begin{Highlighting}[]
\FunctionTok{library}\NormalTok{(garchur)}
\NormalTok{result }\OtherTok{\textless{}{-}} \FunctionTok{garchur}\NormalTok{(y, }\AttributeTok{max\_breaks =} \DecValTok{2}\NormalTok{, }\AttributeTok{trim =} \FloatTok{0.15}\NormalTok{)}
\FunctionTok{summary}\NormalTok{(result)}
\CommentTok{\# Output includes:}
\CommentTok{\# {-} GARCH(1,1) parameter estimates (omega, alpha, beta)}
\CommentTok{\# {-} Unit root t{-}statistic}
\CommentTok{\# {-} Endogenous break dates}
\CommentTok{\# {-} Critical values at 1\%, 5\%, 10\%}
\end{Highlighting}
\end{Shaded}

\subsection{rbfmvar}\label{rbfmvar}

Implements the Residual-Based Fully Modified VAR (RBFM-VAR) estimator of
Chang (2000) for VAR models containing an unknown mixture of I(0), I(1),
and I(2) components. The procedure applies a nonparametric long-run
variance correction to OLS residuals, enabling valid inference without
prior knowledge of integration orders. Supports lag selection by AIC,
BIC, and Hannan-Quinn criteria, modified Wald Granger non-causality
tests, and impulse response functions.

\begin{Shaded}
\begin{Highlighting}[]
\FunctionTok{library}\NormalTok{(rbfmvar)}
\CommentTok{\# Estimate RBFM{-}VAR without specifying integration orders}
\NormalTok{result }\OtherTok{\textless{}{-}} \FunctionTok{rbfmvar}\NormalTok{(}\AttributeTok{data =}\NormalTok{ macro\_data, }\AttributeTok{vars =} \FunctionTok{c}\NormalTok{(}\StringTok{"gdp"}\NormalTok{, }\StringTok{"inflation"}\NormalTok{,}
                  \StringTok{"interest\_rate"}\NormalTok{), }\AttributeTok{lag\_select =} \StringTok{"bic"}\NormalTok{)}
\FunctionTok{summary}\NormalTok{(result)}

\CommentTok{\# Modified Wald Granger causality test}
\NormalTok{granger }\OtherTok{\textless{}{-}} \FunctionTok{rbfm\_granger}\NormalTok{(result, }\AttributeTok{cause =} \StringTok{"inflation"}\NormalTok{,}
                        \AttributeTok{effect =} \StringTok{"gdp"}\NormalTok{)}

\CommentTok{\# Impulse response functions}
\NormalTok{irf }\OtherTok{\textless{}{-}} \FunctionTok{rbfm\_irf}\NormalTok{(result, }\AttributeTok{impulse =} \StringTok{"interest\_rate"}\NormalTok{,}
                \AttributeTok{response =} \StringTok{"gdp"}\NormalTok{, }\AttributeTok{horizon =} \DecValTok{20}\NormalTok{)}
\FunctionTok{plot}\NormalTok{(irf)}
\end{Highlighting}
\end{Shaded}

\subsection{tptest}\label{tptest}

Tests for U-shaped and inverse U-shaped (hump-shaped) nonlinear
relationships in regression models. Extends Lind and Mehlum (2010) to
quadratic, cubic, log-quadratic, and inverse functional forms.
Implements the Sasabuchi (1980) joint slope test, delta-method standard
errors for the turning point, Fieller confidence intervals, and the
Simonsohn (2018) two-lines robustness test. Works as a post-estimation
command after any model returning coefficient vectors and
variance-covariance matrices.

\begin{Shaded}
\begin{Highlighting}[]
\FunctionTok{library}\NormalTok{(tptest)}
\CommentTok{\# Fit a quadratic model}
\NormalTok{model }\OtherTok{\textless{}{-}} \FunctionTok{lm}\NormalTok{(y }\SpecialCharTok{\textasciitilde{}}\NormalTok{ x }\SpecialCharTok{+} \FunctionTok{I}\NormalTok{(x}\SpecialCharTok{\^{}}\DecValTok{2}\NormalTok{), }\AttributeTok{data =}\NormalTok{ ekc\_data)}

\CommentTok{\# Test for U{-}shape or inverse U{-}shape}
\NormalTok{result }\OtherTok{\textless{}{-}} \FunctionTok{tptest}\NormalTok{(model, }\AttributeTok{var =} \StringTok{"x"}\NormalTok{, }\AttributeTok{form =} \StringTok{"quadratic"}\NormalTok{)}
\FunctionTok{summary}\NormalTok{(result)}
\CommentTok{\# Output includes:}
\CommentTok{\# {-} Sasabuchi joint slope test (H0: monotone or inverse U)}
\CommentTok{\# {-} Estimated turning point with delta{-}method SE}
\CommentTok{\# {-} Fieller 95\% confidence interval}
\CommentTok{\# {-} Whether turning point is within data range}

\CommentTok{\# Two{-}lines robustness test}
\NormalTok{robust }\OtherTok{\textless{}{-}} \FunctionTok{tptest}\NormalTok{(model, }\AttributeTok{var =} \StringTok{"x"}\NormalTok{, }\AttributeTok{form =} \StringTok{"quadratic"}\NormalTok{,}
                 \AttributeTok{two\_lines =} \ConstantTok{TRUE}\NormalTok{)}
\end{Highlighting}
\end{Shaded}

\section*{References}\label{references}
\addcontentsline{toc}{section}{References}

\protect\phantomsection\label{refs}
\begin{CSLReferences}{1}{1}
\bibitem[\citeproctext]{ref-Chang2000}
Chang, Yoosoon. 2000. {``Bootstrap Unit Root Tests in Panels with
Cross-Sectional Dependency.''} \emph{Econometric Theory} 16 (6):
1069--74. \url{https://doi.org/10.1017/S0266466600166017}.

\bibitem[\citeproctext]{ref-LindMehlum2010}
Lind, Jo Thori, and Halvor Mehlum. 2010. {``With or Without {U}? The
Appropriate Test for a {U}-Shaped Relationship.''} \emph{Oxford Bulletin
of Economics and Statistics} 72 (1): 109--18.
\url{https://doi.org/10.1111/j.1468-0084.2009.00569.x}.

\bibitem[\citeproctext]{ref-NarayanLiu2015}
Narayan, Paresh Kumar, and Ruipeng Liu. 2015. {``A {GARCH} Model for
Testing Market Efficiency.''} \emph{Journal of International Financial
Markets, Institutions and Money} 41: 121--38.
\url{https://doi.org/10.1016/j.eneco.2014.11.021}.

\bibitem[\citeproctext]{ref-Phillips1995}
Phillips, Peter C. B. 1995. {``Fully Modified Least Squares and Vector
Autoregression.''} \emph{Econometrica} 63 (5): 1023--78.
\url{https://doi.org/10.2307/2171721}.

\bibitem[\citeproctext]{ref-Sasabuchi1980}
Sasabuchi, Syoichi. 1980. {``A Test of a Multivariate Normal Mean with
Composite Hypotheses Determined by Linear Inequalities.''}
\emph{Biometrika} 67 (2): 429--39.
\url{https://doi.org/10.1093/biomet/67.2.429}.

\bibitem[\citeproctext]{ref-Simonsohn2018}
Simonsohn, Uri. 2018. {``Two Lines: A Valid Alternative to the Invalid
Testing of {U}-Shaped Relationships with Quadratic Regressions.''}
\emph{Advances in Methods and Practices in Psychological Science} 1 (4):
538--55. \url{https://doi.org/10.1177/2515245918805755}.

\end{CSLReferences}
